% !PRINTED 709
\weberpage{709}
%<<CONT
on which the determination of the class number still depends, can be
represented by finite expressions. For we have
\[
\frac{1}{n} = \int_0^1 x^{n-1}\,dx,
\]
and accordingly
\[
X = \int_0^1 \Sigma \chi(n)\,x^{n-1}\,dx.
\]

Now $\chi(n)$ remains unchanged when $n$ increases by a multiple of $m$.
But if by $t$ we understand the least positive residue of $n$ and set
\[
n = t + ml,
\]
then $\chi(n) = \chi(t)$, and it follows that:
\[
X = \int_0^1 \overset{t}{\Sigma} \chi(t)\,x^{t-1} \underset{0,\,\infty}{\overset{l}{\Sigma}} x^{ml}\,dx.
\]

If we now set
\weq{2}{f(x) = \overset{t}{\Sigma} \chi(t)\,x^t,}
then $f(x)$ is an integral function of $x$ of degree $m-1$, and at the same
time, in consequence of the relation $\Sigma \chi(t) = 0$:
\weq{3}{f(0) = 0,\quad f(1) = 0,}
and therefore $f(x)$ is divisible by $x(1-x)$. For $X$ there then results,
by the summation formula for the geometric series $\Sigma x^{ml}$:
\weq{4}{X = \int_0^1 \frac{f(x)\,dx}{x(1-x^m)}.}

To find such an integral, the integral calculus prescribes that the
rational fraction under the integral sign be decomposed into partial fractions.
This decomposition into partial fractions gives, however, in consequence of (3), when
\[
r = e^{\frac{2\pi i}{m}}
\]
is set:
\[
\frac{f(x)}{x(1-x^m)} = -\frac{1}{m} \sum\limits_{1,\,m-1}^{s} \frac{f(r^s)}{x-r^s}
\]
[Bd. I, \S.\,29, (11)], and now we have further, by known
%CONT>>
